\section{Executive summary}

This page answers four questions without jargon: what the system does, where it
stands, what is proven, and what has to be decided.

\subsection{What the system does}

LANCE conducts a security audit on a network of connected objects \emph{whose
ownership and authorisation are established}. It explores the network, issues
hypotheses of flaws, attempts to verify them by observable evidence, and then
writes a report. Two policies can carry out this audit: a rule-based automation
and an agent based on a language model. The project seeks to determine which of
the two succeeds better, with identical tools, evidence and budgets.

The audit runs in six phases, each producing a deliverable the next one consumes:
the topology graph, the reconnaissance, the vulnerability hypotheses, the
exploitation attempts, the intrusion campaign, and the report. Every tool call
passes through a single boundary that archives its arguments, its result and its
cost in a journal; that journal, not the policy's own conclusion, is what the
evaluator recounts.

A run can use one of four execution modes. The full pipeline runs all six phases
with a model. A rule-based mode replaces the model with declared rules for
analysis and verification, on phases three and four only. A bare mode gives the
model a bounded tool loop on an address range, without scenario or score. A
scripted mode runs versioned deterministic procedures. Of the four, only the full
pipeline is scored against a ground truth; the others produce results that are
read but not benchmarked.

The system runs in a laboratory: scenarios are deployed on a Proxmox host, audited
over a dedicated network, then torn down. A browser interface, served by three
instances, starts runs, streams their events and reads their results.

\subsection{Where the project stands}

\begin{table}[H]
\centering
\caption{Status in one line per strand.}
\label{tab:etat-volet}
\small
\begin{tabular}{@{}p{3.3cm}p{1.9cm}p{8.3cm}@{}}
\toprule
Strand & Status & Detail \\
\midrule
Audit engine & mature & Six phases in service, four modes, execution boundary that logs every call \\
Evidence & mature & Briefing/oracle/evaluator contract applied to five families; the log is authoritative, not the decision \\
Harness & consolidated & Core shared since 29 October; sandbox executable outside the laboratory \\
Infrastructure & restored & Three instances, automatic delivery pipeline, seven defects fixed \\
Benchmark & partial & One hypothesis measured out of six, one family rejected, no inference computable \\
Scientific proof & absent & No superiority of the model has been established to date \\
\bottomrule
\end{tabular}
\end{table}

\subsection{What is proven, and what is not}

\textbf{Proven.} Every tool call is traced and bounded; a lost piece of evidence
interrupts the run. A harness can run outside the laboratory, under a guard whose
blind spots are declared. On the best-documented family, a third script, which
reads by convention and uses no model, obtains the same result as the agent,
which \emph{absorbs} the gain attributed to the model. The corpus is mostly
detectable by signature, which makes a null gap expected by construction.

\textbf{Not proven.} That a model does better than a competent script. That a gap
is significant. That the results generalise beyond eight cases and one property.

\subsection{The three decisions that govern what follows}

\begin{enumerate}
\item \textbf{The independent families.} The inference tool was repaired on
30 October 2026 and now refuses to produce a number it cannot defend. What
remains is the families themselves: the corpus holds at most three defensible
ones, and a decision needs five. Either build them, applying the filter that one
hour of observation established, or publish a bounded methodological result and
declare it as such.
\item \textbf{The opening condition.} Several families remain undecidable for lack
of a written reference automation. Either write it, or explicitly renounce these
families.

\item \textbf{The comparable arm.} A comparison of the two policies supposes that
both receive the same observations. Four mechanisms currently favour the rules:
the phase-four probes only they may emit, a supplement that rewrites the scan
files, the scanner findings added to every aggregation, and a benchmark score
reserved for the full pipeline. Either give the model an equivalent access, or
state that the comparison measures verification rather than interpretation.
\end{enumerate}

\subsection{How to read this report}

Section~\ref{sec:architecture-pipeline} describes the path of an audit.
Section~\ref{sec:protocoles} describes what makes a result verifiable.
Section~\ref{sec:harnais-sandbox} describes the experimental tooling.
Section~\ref{sec:infrastructure} describes the machine and the delivery.
Section~\ref{sec:benchmark} establishes what is measured.
The conclusion confronts these areas and notes their contradictions; it does not
repeat the summary above.

Each statement refers to the file that supports it. Behaviours that the code does
not allow to be established are flagged as such rather than filled in.

\clearpage
\section{System overview}
\label{sec:vue-ensemble}

Before going into the contracts, a map of the system and of the laboratory. It
situates the components described below and shows where evidence passes.

\begin{figure}[H]
\centering
\begin{tikzpicture}[
  node distance=8mm and 9mm,
  box/.style={draw=line, fill=white, rounded corners=2pt, line width=.55pt,
    text=ink, align=center, inner sep=5pt, font=\sffamily\footnotesize,
    text width=29mm, minimum height=12mm},
  focal/.style={box, draw=accent, fill=rose, line width=.75pt},
  flow/.style={-{Latex[length=1.8mm,width=1.1mm]}, draw=muted, line width=.6pt}
]
\node[box] (lab) {Simulated IoT\\laboratory};
\node[box, right=of lab] (run) {Audit engine\\six phases};
\node[focal, right=of run] (tools) {Execution boundary\\evidence log};
\node[box, below=of lab] (eval) {Evaluator\\offline};
\node[box, right=of eval] (verdict) {Verdict};
\node[box, above=of eval] (art) {Run\\artefacts, 01 to 06};
\node[box, above=of verdict] (oracle) {Private oracle\\fixed before the run};

\draw[flow] (lab) -- (run);
\draw[flow] (run) -- (tools);
\draw[flow] (run) -- (art);
\draw[flow] (art) -- (eval);
\draw[flow] (oracle) -- (verdict);
\draw[flow] (eval) -- (verdict);
\draw[flow, dashed] (tools.south) -- (verdict.north east);
\end{tikzpicture}
\caption{The journey of a piece of evidence. The tools produce logged
observations; the private oracle, written before execution, makes it possible to
recompute the verdict from the log without believing the system's decision.}
\label{fig:tripwire}
\end{figure}

Three properties can be read from this map.

\textbf{The log is the source.} A result is not the conclusion written by the
policy, but what the evaluator recomputes from the archived observations. A cited
piece of evidence that does not appear in the log is rejected.

\textbf{The oracle precedes the run.} It is fixed before the policy runs, and a
finding that contradicts it invalidates the trial instead of correcting it. This
precedence is what prevents the answer from being adjusted after the fact.

\textbf{The sandbox is not the measurement.} It executes the deterministic arms
without a model or a laboratory; it instruments the contrast, it does not measure
it.

\begin{figure}[H]
\centering
\begin{tikzpicture}[
  node distance=3mm and 5mm,
  h/.style={draw=line, fill=white, rounded corners=2pt, line width=.55pt,
    text=ink, align=center, inner sep=3pt, font=\sffamily\footnotesize},
  dev/.style={h, fill=rose, draw=accent},
  flow/.style={-{Latex[length=1.6mm,width=1mm]}, draw=muted, line width=.6pt}
]
\node[h, text width=22mm] (poste) {Development\\workstation};
\node[h, text width=24mm, right=of poste] (ci) {GitHub Actions\\validation};
\node[dev, text width=26mm, right=of ci] (lxc) {Master container\\three instances};
\node[h, text width=22mm, right=of lxc] (umons) {UMONS relay\\model};

\draw[flow] (poste) -- node[above, font=\sffamily\tiny, text=muted] {push} (ci);
\draw[flow] (ci) -- node[above, font=\sffamily\tiny, text=muted] {deploys} (lxc);
\draw[flow] (lxc) -- node[above, font=\sffamily\tiny, text=muted] {calls} (umons);
\end{tikzpicture}
\caption{The delivery pipeline. No execution work takes place on the development
workstation: the push triggers validation, which deploys to the container.}
\label{fig:circuit}
\end{figure}

\clearpage
